% ═══════════════════════════════════════════════════════════════════════════════
% JUSTIFICACIÓN - INSTITUTO ATENEO DE COLIMA
% Estudio Documental: Corrupción en México (2006-2012)
% ═══════════════════════════════════════════════════════════════════════════════

\nota{
  Esta sección articula cinco dimensiones de justificación metodológica: social, 
  académica, práctica, utilidad metodológica y viabilidad, todas fundamentadas en 
  fuentes oficiales verificables y doctrina académica arbitrada.
}

\section{Relevancia social}

La literatura internacional en derechos humanos establece que la corrupción no es un delito sin víctimas, sino una práctica que vulnera directa o indirectamente derechos económicos, sociales y culturales ---entre ellos la vida, la salud, la seguridad y el acceso equitativo a bienes públicos---, afectando de forma desproporcionada a los sectores vulnerables y profundizando la desigualdad estructural \citep{cidh2019,cndh2019,rios2021}. 

En México, el fenómeno opera como un impuesto regresivo de alcance cotidiano: de acuerdo con la Encuesta Nacional de Calidad e Impacto Gubernamental (ENCIG) 2025 del INEGI, publicada en mayo de 2026, 84.1\,\% de la población adulta urbana percibió que los actos de corrupción gubernamental son frecuentes o muy frecuentes, mientras que los pagos indebidos en trámites extrajeron 17,707.4 millones de pesos anuales ---3,865 pesos promedio por víctima---, el costo real más alto de la última década \citep{inegi2026,imco2026}. El análisis retrospectivo de los expedientes clave del periodo 2006--2012 permitió evidenciar cómo decisiones administrativas adoptadas hace más de una década continuaron generando pasivos fiscales, litigios y demandas de justicia vigentes en 2026.

\section{Relevancia académica}

El estudio científico de la corrupción mexicana demanda articular marcos conceptuales internacionales con la tradición académica nacional en ciencia política, sociología, antropología, economía política y derecho público. En esta investigación se integraron aportes clásicos e internacionales ---Nye (1967), Rose-Ackerman (1999, 2004), Garay Salamanca y Salcedo-Albarán (2012), Cancino y el BID (2020), Dávid-Barrett y Fazekas (2020) y Rezzoagli et al. (2021)--- con obras medulares de investigadores mexicanos como Casar y Cejudo (2024), Merino (2024), López Ayllón y Merino (2009), Peschard (2017), Arellano Gault (2012, 2020), Lomnitz (2000), Zaid (2019), Sandoval Ballesteros (2009, 2024), Serna de la Garza (2010), Vázquez Valencia (2019), Astorga (2007, 2015), Flores Pérez (2009) y Reyes Heroles (2003). Dicha articulación permitió contrastar marcos teóricos complejos con resoluciones judiciales y auditorías oficiales verificables.

\section{Relevancia práctica}

Diversos diagnósticos institucionales coinciden en que el principal obstáculo para contener la corrupción en México no radica en la ausencia de legislación, sino en la brecha de implementación entre la norma jurídica y su cumplimiento efectivo \citep{casar2024,lopezayllon2009,peschard2017}. Al reconstruir los factores que hicieron posibles la concentración discrecional del mando policial, la contratación de megainfraestructura sin candados transexenales y la expansión de guarderías subrogadas sin vigilancia independiente, los hallazgos de este trabajo ofrecieron insumos concretos para el diseño de políticas de fiscalización autónoma, digitalización de trámites y control regulatorio sobre servicios concesionados.

\section{Utilidad metodológica}

Desde el punto de vista metodológico, el estudio aportó un procedimiento replicable para investigar fenómenos de gran corrupción mediante un Diseño Anidado Concurrente de Modelo Dominante (DIAC: CUAL + cuan) aplicado a estudios de casos múltiples \citep{hernandez2014}. En lugar de analizar por separado las resoluciones judiciales, las auditorías financieras o el discurso periodístico, se diseñaron y aplicaron dos instrumentos documentales estructurados ---una Ficha de Registro Jurídico-Institucional y Hemerográfico (Apéndice A) y una Matriz Comparativa de Variables, Dimensiones e Indicadores por Caso (Apéndice B)--- que permitieron triangular evidencia probatoria oficial con estadísticas descriptivas secundarias y análisis de encuadre mediático (framing), generando metainferencias trazables.

\section{Viabilidad del estudio}

La investigación resultó plenamente viable desde las dimensiones documental, técnica, temporal y financiera. Se dispuso de acceso abierto a fuentes primarias oficiales y de alta confiabilidad ---informes de la ASF, engroses y dictámenes de la SCJN, recomendaciones de la CNDH, informes de la CIDH, comunicados de la Unidad de Inteligencia Financiera (UIF), expedientes del Departamento de Justicia de Estados Unidos (DOJ), bases estadísticas del INEGI y reportes de Transparencia Internacional---, complementadas con libros y artículos científicos dictaminados. La selección acotada de tres casos principales y dos de contraste contemporáneo garantizó profundidad analítica sin requerir financiamiento externo ni trabajo de campo en entornos de riesgo.